\section{一、问题重述}\label{ux4e00ux95eeux9898ux91cdux8ff0}

\subsection{1.1 问题背景}\label{ux95eeux9898ux80ccux666f}

在现代高端离散制造与电子装备生产中，``多品种、小批量''（High-Mix
Low-Volume,
HMLV）已成为主流生产模式。由于客户定制化需求高、产品生命周期短，制造企业在物料采购与生产排程中面临着高度的不确定性。若按照单纯的静态均值或开环预测安排生产，极易产生巨额的冗余库存导致流动资金严重积压，或在突发需求下频繁发生缺货，损害企业的交付信誉与商业履约率。

因此，如何充分利用海量历史订单数据，科学提取物料需求的多维时序特征，建立兼顾生产提前期时滞、动态安全库存控制、库存资金占用与履约服务水平的闭环滚动排产决策系统，是当前智能制造运筹优化领域的关键核心课题。

\subsection{1.2
题目要求与任务分解}\label{ux9898ux76eeux8981ux6c42ux4e0eux4efbux52a1ux5206ux89e3}

本研究基于某电子制造企业2019年至2022年共22,453条历史订单数据（涉及284种物料），需解决以下四大递进式核心问题：

\begin{enumerate}
\def\labelenumi{\arabic{enumi}.}
\tightlist
\item
  \textbf{问题 1（重点物料筛选与周需求预测模型构建）}：

  \begin{itemize}
  \tightlist
  \item
    依据自然周规则（2019年1月2日所在周为第1周）完成全量历史订单数据的周度聚合清洗；
  \item
    从物料需求出现的频数（下单周数）、数量（总需求量）、趋势与离散度（变异系数\(CV\)）、销售单价以及累计销售额等维度，构建科学综合评价体系，筛选出6种最具代表性的重点物料；
  \item
    建立物料周需求量预测模型，并利用前100周历史数据进行多模型回测与评价。
  \end{itemize}
\item
  \textbf{问题 2（提前期 \(L=1\) 的滚动生产排产与服务水平保障）}：

  \begin{itemize}
  \tightlist
  \item
    生产提前期设定为1周（本周生产物料下周初可用）；
  \item
    从第101周初至第177周初，每周初依据历史信息（\(\le t-1\)周）制定当周排产计划；
  \item
    严格遵循初值条件（第100周末库存和缺货均为零，第100周计划等于第101周实际需求）；
  \item
    确保6种物料在第101～177周的全期平均服务水平（\(1 - \text{缺货量}/\text{实际需求量}\)）不低于85\%，输出典型物料第101～110周明细（表1）及全期综合汇总（表2）。
  \end{itemize}
\item
  \textbf{问题 3（资金占用与服务水平多目标权衡优化）}：

  \begin{itemize}
  \tightlist
  \item
    引入各物料销售单价差异对应的库存资金占用成本率与缺货惩罚损失；
  \item
    调整周生产计划策略，在库存持有资金与履约服务水平间达成Pareto帕累托最优平衡，重新计算并输出表1与表2。
  \end{itemize}
\item
  \textbf{问题 4（提前期 \(L=k \ge 2\) 的通用推广排产）}：

  \begin{itemize}
  \tightlist
  \item
    建立具有多步延迟在途库存的状态转移方程，重新求解提前期\(L=2\)时的问题2与问题3；
  \item
    给出任意提前期\(k\ge 2\)时通用生产计划算法的理论推导、控制律与工程推广方案。
  \end{itemize}
\end{enumerate}
